\begin{textsample}{NEG Dim 2 – vem\_ed\_31 – Score 14.00 – vem\_ed\_31/t161.txt}  \label{ex:f2_neg_177}
VII CBTecLE: desafios e inovações na formação profissional A sétima edição do Congresso Brasileiro de Línguas na Formação Técnica e Tecnológica (CBTecLE) ocorreu em 11 e 12 de setembro de 2025 na Fatec Praia Grande, com o tema “Línguas e Tecnologias na Formação Profissional: desafios e Inovações no Ensino Técnico e Tecnológico”.
Houve 186 inscritos, entre professores e pesquisadores de Etecs, Fatecs, do CPS e de outras 18 instituições de ensino.
Ocorreram 55 apresentações de trabalhos, workshops e pôsteres.
O presidente do CPS, Clóvis Dias, prestigiou a abertura do evento, no Palácio das Artes.
André Luiz Braun Galvão, Coordenador Acadêmico-Pedagógico, representou a Coordenadoria Geral de Ensino Superior de Graduação (CGESG / CPS).
A mesa de abertura contou ainda com a presença de Carla Aparecida Pedriali Moraes, chefe da Divisão de Ensino e Pesquisa no Ensino Superior (Depes / CGESG / CPS); Ulysses Diegues, coordenador da Fatec Praia Grande, e autoridades da prefeitura do município.
Após homenagens a Mariane Teixeira, idealizadora do CBTecLE e membro da comissão organizadora das seis primeiras edições do congresso, a Elenir Silva, membro da comissão geral do VII CBTecLE, e a Osvaldo Succi Junior, coordenador da equipe de Apoio à Internacionalização do Ensino Superior da Depes / CGESG / CPS, teve início a primeira mesa-redonda, com o tema “Os desafios do ensino de LinFE na contemporaneidade”.
Antonio Ferreira da Silva Júnior (UFRJ), Catia Veneziano Pitombeira (UFAL) e Marcela Iochem Valente (UERJ) debateram conceitos e fundamentos teóricos relativos ao ensino de Línguas para Fins Específicos (LinFE).
Na tarde de 11 de setembro, a equipe de Apoio à Internacionalização do Ensino Superior da Depes / CGESG / CPS apresentou o workshop"Projetos Colaborativos Internacionais na era da IA: ferramentas para análise".
Patrícia Patrício e Neusa Haruka Sezaki Gritti fizeram uma breve introdução sobre os PCIs e suas etapas (desde a busca de parceiros até a avaliação dos projetos).
Regiane Camargo mostrou ferramentas de IA para planejamento de atividades na plataforma Padlet e Maria Claudia Delfino demonstrou o uso da ferramenta Notebook LM com o corpus dos 30 números de VEm.
O conjunto das publicações serviu como referencial para levantar perguntas sobre Intercâmbios Virtuais, elaborar um mapa mental, um podcast e até um vídeo de apresentação resumindo os conteúdos das edições.
Nos dois dias do CBTecLE, professores de Fatecs que realizaram PCIs compartilharam suas experiências nas sessões de comunicação oral na Fatec Praia Grande: “Entre saberes, linguagens e identidades: a experiência de um PCI no currículo do ensino superior” foi o relato de José Carlos Barbosa Lopes sobre duas edições (segundo semestre de 2020 e primeiro de 2021) do projeto desenvolvido nas Fatecs Ipiranga e São Paulo com Jamestown Community College (JCC SUNY / EUA). “Internacionalização do currículo na prática: um relato de Projeto Colaborativo Internacional entre Brasil e EUA”, pelas professoras Kátia Cristina Galatti e Talita Botelho Nunes, da Fatec Taquaritinga.
Elas refletiram sobre o PCI realizado com Perimeter College / Georgia State University (EUA), uma proposta interdisciplinar envolvendo Língua Inglesa, Agroecologia e Biologia e alinhada com os Objetivos de Desenvolvimento Sustentável (ODS).
O projeto foi desenvolvido em duas etapas: no segundo semestre de 2024, os estudantes trabalharam o tema “Hortas sustentáveis”, pesquisando técnicas e oferecendo os produtos da horta a uma instituição de caridade.
No semestre seguinte, as equipes compararam esforços de conservação ambiental nos estados da Geórgia e de São Paulo. “A discussão foi enriquecida por dados locais e entrevistas com membros das comunidades”, ressaltam as autoras. (Leia mais sobre o PCI no Artigo de Opinião). “Educação sem fronteiras: o papel dos Projetos Colaborativos Interna- cionais (PCIs) na formação e empregabilidade dos alunos da Fatec de Bragança Paulista”, de Natalie Goes e colegas do Núcleo de Estudos da Linguagem (NELF) da unidade.
O grupo publicou Artigo de Opinião em VEm 28 detalhando essa proposta: http://bit.ly/4qQdMWP “Escrituras Narrativas: cinco edições de um projeto internacional de narrativas audiovisuais no ensino de espanhol”, de Odenildo França Almeida.
Reflexões sobre o amadurecimento do PCI desenvolvido entre Fatec Ipiranga e Uniminuto (Colômbia) foram compartilhadas pelo autor, que também publicou artigo sobre o projeto em VEm 28.
Maria Edna da Silva Gomes (Fatec Barueri) apresentou o pôster"Espanhol para fins específicos e formação crítica: sentidos construídos por alunos de Comércio Exterior em Projeto Internacional Colaborativo".
A autora publicou Artigo de Opinião sobre o PCI em VEm 26: https://bit.ly/47G7v8 g Na manhã do dia 12, houve a segunda mesa-redonda, intitulada"Inteligência Artificial e Educação na Atualidade"e composta por Tony Berber Sardinha (PUC-SP) e Mariana Reis Mendes (PUC-SP, UFPEL, UFRGS).
A mediação ficou a cargo de Fernanda Peixoto Coelho (CPS).
Mariana citou as recomendações da Unesco para o uso de IA: mentalidade centrada no humano, aprender com, para e sobre IA.
Defendeu uma postura responsiva e responsável, com educação midiática e multiletramento engajado.
E apresentou algumas ferramentas para educadores, como o Projeto Inclu.ai para Design Universal para Aprendizagem, o Notebook LM e o app Gemini para conversação e prática argumentativa.
Tony Berber Sardinha alertou para o fato de que as IAs Generativas produzem textos por meio de probabilística e a partir de bases de dados da internet (blogs, sites de notícias, artigos acadêmicos).
As IAs Generativas “não têm conhecimento” do que é um texto, muito menos dos diversos gêneros textuais. À tarde, Adolfo Tanzi Neto (UFRJ) ministrou a conferência “Mediação, Agência e Artefatos em Contextos de Ensino-Aprendizagem: uma abordagem sociocultural para transformação social em LinFE”.

% matched lemmas: 
\end{textsample}
